\section{Introduction}\label{sec:objective}
%%% ====================================================================

What determines the values $\zeta(3)$, $\zeta(5)$, $\zeta(7)$, \ldots?
Euler showed that the even values $\zeta(2k) = r_k\pi^{2k}$
($r_k \in \mathbb{Q}$) are governed by $\pi$ alone.
The odd values have resisted a comparable structural account for
over two centuries: Ap\'ery~\cite{Apery} proved
$\zeta(3)\notin\mathbb{Q}$,
Ball and Rivoal~\cite{BallRivoal} showed that infinitely many
$\zeta(2k+1)$ are irrational, but no organising principle has
emerged.  This paper identifies such a principle: the pentagonal
arithmetic of $\mathbb{Q}(\sqrt{5})$, generated by
$\golden^2 = \golden + 1$.  The irrationality itself traces to
$\log\golden$ being transcendental (Lindemann--Weierstrass).

The key identity is the Dedekind factorisation
\begin{equation}\label{eq:main-factorization}
  \zeta(s) = \frac{\ZKsqrt(s)}{L(s,\chi_5)},
\end{equation}
where the Artin character $\chi_5$ is dictated by the pentagon:
\begin{equation}\label{eq:chi5-pentagon}
  \chi_5(a) = \frac{\log|2\cos(\pi a/5)|}{\log\golden}\in\{+1,-1\}.
\end{equation}
This character classifies every prime $p$ as split or inert in
$\mathbb{Z}[\golden]$, determining the local Euler factor $f_p(s)$
and hence the entire product $\ZKsqrt(s) = \prod_p f_p(s)$.
The base case $L(1,\chi_5) = 2\log\golden/\sqrt{5} = T/(\pi\sqrt{5})$
($T = 2\pi\log\golden$) shows that $\pi$---the circumference of the
circle in $T^2_{\mathrm{PCF}}$---is present from the outset.
The result resolves the apparent freedom noted by
Elvang et al.~\cite{Elvang2026}: the Wilson coefficients
$a_{2k-1,0} = \alpha^{\prime\,2k+1}\zeta(2k+1)$ that are free
in the $\mathcal{N}{=}4$ SYM EFT are fixed by this arithmetic.

The complete chain traced in this paper is:
\[
\begin{tikzcd}[column sep=1.2em, row sep=2.5em]
  & & & {\textstyle\prod f_p} \ar[r] & \zeta_K \ar[r]
    & \zeta(2k{+}1) \\
  \golden^2{=}\golden{+}1 \ar[r]
    & \psi_p \ar[r]
    & \chi_5 \ar[r]
    & G \ar[u] \ar[d] & & \\
  & & & \mu_3{=}\tfrac{1}{2} \ar[r]
    & S_3 \ar[r]
    & \omega^k/2
\end{tikzcd}
\]
Both branches originate from the same Galois functor
$G\colon\Ztwenty\to\mathbb{Z}_2\times\mathbb{Z}_2$.
The upper branch produces the odd zeta values: $G$ classifies
primes, determining every $f_p$, assembling the Euler product,
and delivering $\zeta(2k+1) = \ZKsqrt(2k+1)/L(2k+1,\chi_5)$
(Sections~\ref{sec:frobenius-split}--\ref{sec:pentagon};
worked example for $\zeta(3)$ in~§\ref{ssec:zeta3-example}).
The lower branch produces the spectral structure: $|\ker G| = 2$
forces $\mu_3 = 1/2$, which determines $S_3$ symmetry, eigenvalues
$\omega^k/2$, the AdS radius $\ell = 1$ (from S-duality at
$\tau_{\mathrm{PCF}} = i$), and the Brown--Henneaux~\cite{BrownHenneaux}
central charge $c = 3\ell/(2G_N) = 3$
(Section~\ref{ssec:unifying}, Lemmas~\ref{lem:spectral-re}
and~\ref{lem:ads-radius}).
The regulator $R = \log\golden$ links both branches: it is the
arithmetic size of the field ($L(1,\chi_5) = 2R/\sqrt{5}$), the
worldsheet period ($T = 2\pi R$), and the entropy anchor
($S_{\mathrm{BH}}/k_B = \lambda_{\log}\cdot R$).
Section~\ref{sec:categorical-framework} makes the categorical content
explicit: $\{\psi_p\}$ is a functor, the Euler product is a colimit,
the Dedekind factorisation is a natural isomorphism across the four
quadratic fields classified by~$G$, and the two branches form a span
with common source~$G$.

This span discharges the \texttt{bridge\_axiom} of the Primitive
Structure~\cite{V11} (\emph{The Hilbert P\'olya Operator and the
Primitive Structure of the Complex Plane}, hereafter
\emph{the Primitive Structure}): because~$G$ simultaneously
produces~$\zeta(s)$ (via the Euler product~\cite{Euler} and analytic
continuation~\cite{Riemann}) and the spectral constraint
$\mu_3 = 1/2$ (via $|\ker G| = 2$), the same map that
determines~$\zeta(s)$ everywhere also constrains
$\mathrm{Re}(\rho)$.
The Primitive Structure establishes the functorial transport of the
eigenvalue constraint $\mathrm{Re}>0 \Rightarrow \mathrm{Re}=1/2$
from $\OmH$ to~$\zeta$ across the full critical strip; the present
paper provides the \emph{mechanism} underlying that transport.
Theorem~\ref{thm:unifying} shows that every structure in this
paper derives from the single relation $\golden^2 = \golden + 1$.

%%% ====================================================================
\section{Context}\label{sec:context}
%%% ====================================================================

In the Primitive Structure~\cite{V11}, the proof that
$\mathrm{Re}(\rho)=1/2$ within the PCF categorical framework proceeds
through a squeeze (Theorem~10 of~\cite{V11}): two independent bounds, both derived from the
arithmetic of $\Ztwenty$ and the Galois functor
$G\colon\Ztwenty\to\mathbb{Z}_2\times\mathbb{Z}_2$
(Definition~5 of~\cite{V11}), force equality
at~$1/2$.  The squeeze itself is verified in Lean~4 with
$0$~\texttt{sorry}.  However, the step connecting the categorical
machinery to the zeros of~$\zeta$---the identification of
$\mathrm{Re}(\rho)$ with an eigenvalue of~$\OmH$---is left as an
axiom (\texttt{bridge\_axiom}, equivalently the structure field
\texttt{bounded\_by\_mu :\ $\rho_{\mathrm{re}}\le\mu_3$} of
\texttt{TernaryLZero}, Definition~7 of~\cite{V11}).

The axiom states $\rho_{\mathrm{re}}\le\mu_3=1/2$: it is
Premise~A of the squeeze (Proposition~17 of~\cite{V11}).
Once provided, \texttt{hecke\_1920}
derives the companion bound $1-\rho_{\mathrm{re}}\le\mu_3$
(Premise~B, Proposition~18 of~\cite{V11}), and the two together force
$\rho_{\mathrm{re}}=1/2$ by \texttt{linarith}.  The mathematical
justification is the Spectral Embedding (Proposition~15
of~\cite{V11}): the Euler product enters~$\Ztwenty$, the Galois
functor~$G$ classifies primes into four characters, the $S_3$
symmetry of the tripartite object restricts
$\mathrm{spec}(\OmH)=\tfrac{1}{2}\{1,\omega,\omega^2\}$ with
real parts $\{1/2,\,-1/4\}$, both ${\le}\,\mu_3$.  In Lean the
step remains an axiom.
The Primitive Structure
cites the Euler product identity, Gelfand's Banach algebra
theory~\cite{Gelfand}, and the Dunford--Schwartz spectral mapping
theorem~\cite{DunfordSchwartz}
as justification, but does not demonstrate that these three
classical tools act on a common structure: the Euler product
organised by~$G$ through the same map that produces
$\mathrm{spec}(\OmH)$.

The present paper develops this bridge.  The development formalises
the relationships between the $\mathbb{F}_1$-geometric programme
(Borger~\cite{Borger}, Manin~\cite{Manin},
L\'opez Pe\~na--Lorscheid~\cite{Lorscheid}) and the PCF framework
in relation to the Euler product:
the $\Lambda$-ring structure of~$\Rpcf$---the Frobenius lifts
$\psi_p(\golden)=\golden^p$ with composition law
$\psi_p\circ\psi_q=\psi_{pq}$---is simultaneously the
$\mathbb{F}_1$-descent data in the sense of Borger and the source of
every local Euler factor $f_p(s)$ through the Fibonacci splitting
criterion $\chi_5(p)\equiv F_p\pmod{p}$
(Proposition~\ref{prop:fib-split}), with the classification of
residue classes in $\Ztwenty$ via the golden ratio due to
Grisales Herrera~\cite{GrisalesHerrera}.  Establishing the categorical
infrastructure---the Frobenius system as a functor
(\S\ref{ssec:frobenius-functor}), the Euler product as a colimit
(\S\ref{ssec:euler-colimit}), the Dedekind factorisation as a natural
isomorphism (\S\ref{ssec:dedekind-natural}), and the arithmetic--spectral
span (\S\ref{ssec:span})---we found that the same chain that discharges
the bridge also yields a structural determination of the odd zeta values:
$\zeta(2k+1) = \ZKsqrt(2k+1)/L(2k+1,\chi_5)$, with every component
fixed by the pentagonal arithmetic of $\mathbb{Q}(\sqrt{5})$ and no
free parameter.  This determination connects to a problem even older
than the Riemann Hypothesis---the structural origin of the odd zeta
values---and relates it to string theory from a viewpoint that both
complements and constrains recent independent developments.

Elvang, Herderschee and Morales~\cite{Elvang2026} demonstrate that in the
weakly-coupled 4d planar $\mathcal{N}=4$ SYM EFT, the Wilson coefficients
\[
  a_{2k-1,0} = \alpha^{\prime\,2k+1}\zeta(2k+1), \qquad k = 1, 2, 3, \ldots
\]
are free parameters: supersymmetry and tree-level factorisation fix all
$a_{2k,q}$ in terms of lower-order coefficients but leave $a_{2k-1,0}$
unconstrained.  In the Veneziano amplitude~\cite{Veneziano} these take the values
$\alpha^{\prime\,2k+1}\zeta(2k+1)$, but from the EFT alone the odd zeta
values enter as external data.

The PCF (Precedent-Current-Forthcoming) framework,
formalised in the Primitive Structure~\cite{V11}, provides the missing arithmetic structure.
The present paper shows that once the Frobenius action of $\golden$ on the
primes of $\mathbb{Q}(\sqrt{5})$ is taken into account, the odd zeta values
are determined---not through new closed forms, but through the structural
identity~\eqref{eq:main-factorization} together with the class-number formula
for $\mathbb{Q}(\sqrt{5})$.

%%% ====================================================================
